\section{外交关系法和领事关系法}
\begin{enumerate}
  \item 外交关系以设使馆为标志，领事关系以设领馆为标志。
  \item 区别。
        \begin{enumerate}
          \item 使馆。中央机构；全局；全境；首都；一个。
          \item 领馆。相应的地方机关；商务文化侨民保护，\zd{没有政治职能}；辖区；不限。
          \item 二者\zd{平等}，各司其职，没有隶属关系。
        \end{enumerate}
  \item 特别使团。临时派出。
  \item 使馆。独立的法律地位，能够自行行使权利，承担义务。
        \begin{kk}{级别}
          大使馆、公使馆、代办处 \yjt 大使、公使、代办。\zd{职务相同，礼遇不同}。
          \begin{enumerate}
            \item 大使和公使。国家元首出面迎接，最高可约见到国家元首。
            \item 代办。国外交部长出面，最高可约见到外交部长。
          \end{enumerate}
        \end{kk}

        \begin{kk}{三大类人员}
          \begin{enumerate}
            \item 外交人员。馆长、参赞、武官、外交秘书、随员 \yjt \zd{不受欢迎}的人（无需理由）。
            \item 行政人员和服务人员 \yjt \zd{不被接受}的人。特权和豁免仅限\zd{职务范围内}。
          \end{enumerate}
          \tcblower
          可以在任何时候，不说明理由的宣布。

          没来不给来；来了 \yjt\, Not(酌情召回/就地免职) \yjt 不承认 \yjt 限期离境。
        \end{kk}
  \item 派遣和任职先行向接受国提名。\zd{长官特使派外人可撤回}。
  \item 职务开始。使馆馆长\zd{交}国书，领馆馆长\zd{接}过领事证书到其辖区，其他人员到达职务自动开始。
  \item 馆舍的特权与豁免。总体规律：\zd{使馆特权比领馆高}。
        \begin{enumerate}
          \item 使馆馆舍。\qd{扩大解释，公务区休息区私人官邸}，\zd{绝对不得侵犯}，馆长同意。
          \item 领馆仅限工作区域，紧急情况推定同意，馆舍设备可被征用事后补偿。
          \item 通讯自由，装置及使用无线电\zd{发报机}均须许可。
          \item 外交信差和领事信差\zd{执行职务时}人身不得侵犯。
          \item 可交\zd{商业机长}转递但\zd{不是为外交信差身份}。
          \item 使馆外交邮袋绝对不得侵犯权；领馆\zd{重大理由，截下来，当面 \hz 退回}。
        \end{enumerate}
  \item 人员的特权与豁免。外交人员、领事官员\zd{执行职务行为}。
        \begin{enumerate}
          \item \zd{豁免的是管辖，不是刑事责任}，外交途径而非司法途径。
          \item 外交人员绝对人身，领事官员\zd{严重罪行 \hz 执行已生效判决}。
          \item \zd{免直接税不免费}。
          \item 没有作证义务，\zd{不能个人放弃，派遣国明示放弃}。
          \item 外交人员范围扩大配偶\zd{未成年子女}，取得中国国籍/定居资格则缩小到职务范围。
          \item \qd{特别使团性质上身份上是外交人员，特权与豁免上 \ydy 领馆领事员}。
          \item 时间。入境至离境/离境合理期间享有。
        \end{enumerate}

        \begin{kk}{管辖问题}
          \begin{enumerate}
            \item 先看身份，4 种：私有不动产、继承、公务外专业/商事纠纷、主动起诉、反诉；领事官员看职务。
            \item 是否启动\zd{内部报告程序}。主体为\zd{被告、第三人，决定受理之前}。
          \end{enumerate}
        \end{kk}
\end{enumerate}
